%% ma: L12-B08-0089
%% lop: 12
%% chuong: 2
%% bai: 8
%% dang: 12.8.12
%% loai: TN4
%% muc: TH
%% dap_an: "A"
%% nguon: "CĐ2.2.pdf giáo viên gửi 27/10/2026 (M07, Lê Minh Khôi – Nguyễn Anh Khôi, Chương II Chủ đề 2 Tọa độ của vectơ trong không gian), Câu 23"
%% kien_thuc: "Bài 8 (tọa độ vectơ, hai vectơ bằng nhau)"
%% trang_thai: chinh-thuc
\begin{ex}
Trong không gian $Oxyz$, cho ba điểm $A(1;2;-1)$, $B(2;3;2)$ và $C(0;1;4)$. Tọa độ của điểm $D$ để tứ giác $ABCD$ là hình bình hành là
\choice{\True $D(-1;0;1)$}{$D(1;0;1)$}{$D(3;4;5)$}{$D(-1;2;1)$}
\loigiai{$ABCD$ là hình bình hành $\Leftrightarrow\overrightarrow{AD}=\overrightarrow{BC}=(-2;-2;2)\Rightarrow D(-1;0;1)$ (ba điểm $A$, $B$, $C$ không thẳng hàng).}
\end{ex}
